\subsection{Finite public-carrier census}

To prevent repeated searches in non-identifying datasets, we froze eight
public candidates and applied five gates: same carrier, an $x_Q$-capable state
distance, an action-like observable, an independently once-counted $x_A$, and
errors plus recovery information.  Table~\ref{tab:public-census} summarizes
the result.  ``Partial'' means that the required observable exists only as a
model-dependent or same-pipeline proxy.

\begin{table}[ht]
\centering
\small
\caption{Public-carrier census. No candidate passes the independent-$x_A$ gate.}
\label{tab:public-census}
\begin{tabularx}{\textwidth}{@{}Xccccc@{}}
\toprule
Carrier/data & same & $x_Q$ & action-like & independent $x_A$ & errors/recovery\\
\midrule
Tan superconducting qubit & yes & yes & yes & no & partial\\
Yu NV centre & yes & yes & no & no & partial\\
Kim black phosphorus & yes & yes & yes & no & partial\\
Zheng non-Abelian circuit & yes & yes & yes & no & partial\\
Pires NMR open QSL & yes & yes & yes & no & partial\\
Ness trapped atom QSL & yes & yes & yes & no & yes\\
KIT mitigated tomography & yes & yes & no & no & yes\\
Photonic overlap data & yes & yes & no & no & partial\\
\bottomrule
\end{tabularx}
\end{table}

The strongest standard-QOR control is the trapped-atom dataset of Ness et
al.\ \citep{Ness2021QSL}: it publishes two-time overlaps together with mean
energy, energy variance, timing and an associated data archive.  It can test
the ordinary speed-limit side, but energy moments are not an independently
defined morphological action separation.  We audited all 53 published
Figure~2 overlap points against the larger of the interpolated
Mandelstam--Tamm and Margolus--Levitin lower limits.  Pointwise central values
occasionally fall below the fitted limit, but after combining the measured
lower and bound upper uncertainties in quadrature the largest one-sided
deviation is only $1.4465\sigma$; no point exceeds $1.96\sigma$.  Thus the
public control is consistent with the standard QSL and supplies no anomalous
Tau residual.

The strongest spectral-response
candidate is the public black-phosphorus package of Kim et al.\
\citep{Kim2025QGTData}, which contains ARPES dispersion and QGT products.  Its
$x_A$ map is not source-frozen, and the QGT and candidate spectral response
share ARPES/pseudospin processing.  The non-Abelian superconducting-circuit
\citep{Zheng2022QGT} and open-NMR QSL \citep{Pires2024QSL} candidates fail for
the analogous same-pipeline reason.

There is also a structural obstruction to simply declaring spectral
curvature to be $x_A$.  For the generic two-band family
\begin{equation}
 H(q)=r(q)\,\mathbf n(q)\!\cdot\!\boldsymbol\sigma,
 \qquad
 r(q)=1+bq^2,
 \qquad
 \mathbf n(q)=(\sin aq,0,\cos aq),
\end{equation}
the ground-state quantum metric and positive-band curvature at $q=0$ are
\begin{equation}
 g_{qq}(0)=\frac{a^2}{4},
 \qquad
 E_+''(0)=2b.
\end{equation}
Holding $a$ fixed while varying $b$ gives the same QGT and different spectral
curvature; holding $b$ fixed while varying $a$ gives the same spectrum and
different QGT.  Hence neither quantity determines the other without an
additional source-owned law tying radial spectral data to angular state data.
This countermodel rules out relabelling an ARPES band curvature as an
independent morphological action record.

The same obstruction survives when a public calculation supplies separate
energy- and fidelity-susceptibility estimators, as in the open PMR-QMC packet
of Ezzell et al.\ \citep{Ezzell2024Susceptibility}.  For a nondegenerate
ground state of $H(\lambda)=H_0+\lambda V$, standard perturbation theory gives
\begin{equation}
 \chi_E=-\partial_\lambda^2E_0
 =2\sum_{n>0}\frac{|V_{n0}|^2}{\Delta_n},
 \qquad
 \chi_F=\sum_{n>0}\frac{|V_{n0}|^2}{\Delta_n^2}.
 \label{eq:susceptibility-moment-nogo}
\end{equation}
They are different spectral moments.  Already for one excited level,
$\chi_E/\chi_F=2\Delta$, so no parameter-independent universal scale converts
one into the other.  Independent estimation is therefore necessary but not
sufficient: identifying an action-side $x_A$ requires a source-owned rule that
fixes the gap weighting and normalization without using the fidelity data.
The public PMR-QMC package is useful as a computational control of the two
moments, not as an occupied physical carrier or a test of
Eq.~\eqref{eq:tau-discriminator}.

Consequently the census has zero full passes.  Ness 2021 is retained only as
a standard QOR/recovery control.  Kim 2025 may be reopened only after an
action/work/response-curvature map is defined without using the state or QGT
data that produce $x_Q$.  This is the stopping verdict for the present public
inventory; adding another tomography, metric or Berry estimator cannot
reduce the finite discriminator blocker.  The complete Ness tables, provenance,
uncertainty audit and two-band countermodel are included in the reproducibility
package.

\subsection{Operational boundary for the action-side observable}

The preceding audits permit a finite exhaustion statement for the standard
public-observable routes.  Let $\mathfrak O_Q$ be the algebra generated by
state tomography, Hamiltonian spectroscopy, quantum-geometric response and
QOR timing on the selected carrier.  Terminally equivalent completions can
agree on all of $\mathfrak O_Q$ while assigning
\begin{equation}
 x_A^{(0)}=x_Q,
 \qquad
 x_A^{(\lambda)}=x_Q+\lambda x_Q^2,
 \qquad \lambda>0.
 \label{eq:terminal-equivalent-actions}
\end{equation}
Therefore no statistic measurable only in $\mathfrak O_Q$ identifies finite
A8b.  The concrete candidate reductions fail for complementary reasons:
\begin{enumerate}[label=(O\arabic*)]
\item QGT, fidelity and tomography are functions of $\mathfrak O_Q$ and are
  circular as action-side estimators;
\item reversible quasistatic work vanishes for a fixed-gap rotating ground
  state although its endpoint Uhlmann chord is nonzero;
\item band curvature is independent of the eigenstate QGT in the two-band
  countermodel above;
\item energy susceptibility and fidelity susceptibility are inequivalent
  spectral moments by Eq.~\eqref{eq:susceptibility-moment-nogo}.
\end{enumerate}

\paragraph{Operational $x_A$ no-go.}
Within the frozen public inventory and the observable algebra
$\mathfrak O_Q$, there is no operationally identified independent estimator
of $x_A=\mathcal A_X^{PC}/A_*$.  The common-Gram equality remains a conditional
cross-domain prediction, not an executable public-data test.

This no-go does not make the prediction unfalsifiable in principle.  It fixes
the minimum admissible reopening packet.  One must independently measure a
source/body-side local stiffness $K_A(q)$, freeze its action unit and
continuation law before examining withheld quantum endpoints, and then test
the blind predictions
\begin{equation}
 \widehat x_A(q_i,q_j)
 \stackrel{\rm blind}{=}
 2\bigl(1-\sqrt{F(\rho_i,\rho_j)}\bigr)
 \label{eq:blind-action-prediction}
\end{equation}
on non-collinear pairs and closed triangles.  Triangle inconsistency would
detect unmodelled holonomy or finite-response corrections.  Relative
distance ratios can reject nonlinear laws before $A_*$ is known, but full
A8b additionally requires the common absolute action unit.  No located public
dataset contains this independent stiffness-plus-withheld-endpoint packet.

Local stiffness alone is still insufficient.  Consider two amplitude-plane
embeddings
\begin{equation}
 W_0(q)=(q,0),
 \qquad
 W_\kappa(q)=\frac{1}{\kappa}
 \bigl(\sin\kappa q,1-\cos\kappa q\bigr).
 \label{eq:isometric-curve-countermodel}
\end{equation}
Both have the identical local pullback metric
$\|\partial_qW\|^2=1$, but for an endpoint interval $\Delta q$ their chords
are
\begin{equation}
 d_0=|\Delta q|,
 \qquad
 d_\kappa=\frac{2}{|\kappa|}
 \left|\sin\frac{\kappa\Delta q}{2}\right|,
\end{equation}
which differ at finite separation.  Thus $K_A$ determines only intrinsic
infinitesimal length.  A blind finite certificate must additionally freeze
the connector's extrinsic continuation, or an equivalent amplitude lift and
connection, and test its holonomy.  The minimal reopening packet is therefore
\begin{equation}
 \boxed{\{K_A,\ C_q,\ \nabla^C,\ A_*,\ \text{withheld endpoints and loops}\}}.
 \label{eq:minimal-blind-packet}
\end{equation}
On a simply connected tested domain, flat right-unitary-basic transport makes
the integrated lift path independent and permits
Eq.~\eqref{eq:blind-action-prediction}.  Nontrivial descriptor-visible
holonomy must be retained as physical path data rather than quotiented away.

Even all pairwise endpoint distances do not determine that loop datum.  For
the three pure qubit states
\begin{equation}
 |\psi_0\rangle=(1,0),\quad
 |\psi_1\rangle=(\cos\theta,\sin\theta),\quad
 |\psi_2^{\pm}\rangle=(\cos\theta,e^{\pm i\phi}\sin\theta),
\end{equation}
the $+$ and $-$ packets have identical three pairwise fidelities.  Their
gauge-invariant Bargmann products
\begin{equation}
 \mathcal B_\pm=
 \langle\psi_0|\psi_1\rangle
 \langle\psi_1|\psi_2^\pm\rangle
 \langle\psi_2^\pm|\psi_0\rangle
\end{equation}
are complex conjugates and carry opposite nonzero phases.  Pairwise chords
therefore cannot certify the connector's phase transport.  At least one
closed triangular loop witness is mandatory in addition to non-collinear
endpoint tests.  For mixed states the corresponding object is the declared
Uhlmann transport/holonomy, not an arbitrary purification phase.

\subsection{Alternative physical sources}

The operational no-go above excludes reconstruction of $x_A$ from the frozen
standard quantum-observable algebra; it does not exclude a different physical
measurement source. We therefore performed a finite same-carrier audit of
six experiments that realize complementary parts of the blind packet.
Batalh\~ao et al. directly reconstructed a quantum work distribution in an
NMR spin system \citep{Batalhao2014Work}. Cottet et al. combined microwave
work extraction with full state tomography in a superconducting quantum
Maxwell-demon experiment \citep{Cottet2017Demon}. Mixed-state geometric phase
was measured interferometrically in NMR \citep{Ghosh2006MixedPhase}, and an
Uhlmann phase was implemented with superconducting qubits
\citep{Viyuela2018Uhlmann}. Finally, an Xmon experiment combined a closed
geometric gate with process tomography \citep{Zhao2020GeometricGate}.
Naghiloo et al. additionally tracked individual transmon trajectories and
assigned heat and work along the same reconstructed state path
\citep{Naghiloo2020Trajectories}.

These results establish the separate experimental availability of all three
measurement legs,
\begin{equation}
 \text{work/action response},\qquad
 \text{endpoint state geometry},\qquad
 \text{closed-loop phase or holonomy}.
 \label{eq:alternative-source-legs}
\end{equation}
They do not yet form one source-frozen certificate. No audited publication
simultaneously supplies all three legs on the same physical carrier, an
independent map from the measured work-like observable to the body action,
and a reusable numerical archive sufficient for the full test. In
particular, thermodynamic work and imposed control-pulse area are standard
terminal observables; neither may be renamed
$\mathcal A_X^{PC}$ without a source-owned calibration theorem.
The Naghiloo packet is the strongest trajectory-resolved candidate, but its
work is computed from monitored $\rho(t)$ and known $H(t)$; it therefore adds
no independent action coordinate, reports no closed Bargmann/Uhlmann witness,
and has no located reusable raw-trajectory archive.
Its frozen publication-level readiness score is therefore $6/13$. The reported
$20\,$ns sampling, approximately $10^4$ trajectories, measurement efficiency
$0.35$, heat--feedback correlation $-0.68$, and feedback/open-loop regression
slopes $0.437/0.006$ support strong first-law and feedback controls. They do
not license the finite A8b, holonomy or QOR scores.

The Cottet packet plays a complementary and stronger instrumental role.
Extracted work is measured directly as an input--output microwave-power
difference and checked against an independent qubit-energy measurement and
full cavity-memory tomography. Its readiness score is $7/13$. This is an
independent measurement instrument, not yet an information-independent
terminal coordinate. Reported controls include informed/ignorant memory
occupations $\bar n=9/0$, residual qubit excitation $0.027\pm0.010$, and memory
entropies $1.0\pm0.05$ versus $1.2\pm0.1$. It still does not supply a
source-frozen map from this terminal work to $\mathcal A_X^{PC}/A_*$ or a
closed holonomy. It is therefore the primary directly measured energy control,
while the Naghiloo packet remains the trajectory control; their records may
not be stitched into one Tau certificate.

The distinction follows directly from input--output theory. With frozen input
field and coupling controls $c$, the measured output satisfies
\begin{equation}
 b_{\rm out}=b_{\rm in}-\sqrt{\gamma_b}\,\sigma_- ,
 \qquad
 \overline P_{\rm out}=F(\rho_S,c).
 \label{eq:cottet-factorization}
\end{equation}
Consequently
$D\overline P_{\rm out}=D_\rho F\,D\rho_S$ and stacking mean output power with
complete qubit tomography does not increase the conditional descriptor rank.
The power detector is physically independent and provides an excellent
cross-check, but its mean record is not information-novel once $(\rho_S,c)$ is
known.

There is a precise qualification for fluctuations. Multi-time output moments
need not factor through an instantaneous $\rho_S(t)$ alone: they also depend
on the dynamical generator $\mathcal L$, coupling operators, input-field state
and controls. Under a frozen Markov input--output model, however, quantum
regression gives schematically
\begin{equation}
 \langle O_n(t_n)\cdots O_1(t_1)\rangle
 =\mathcal F_n(\rho_S(0),\mathcal L,L_b,\rho_{\rm in},c).
 \label{eq:cottet-correlation-factorization}
\end{equation}
Correlations may therefore be novel relative to endpoint or instantaneous
tomography while remaining redundant relative to independently frozen process
tomography. A residual surviving the latter test diagnoses an incomplete
standard process model, not automatically the Tau body action. The Cottet
publication-level packet reports averaged power and no reusable single-shot
archive sufficient for this residual test, so no fluctuation score is licensed.

We therefore performed a finite direct-public-data audit. The Wang
longitudinal-readout Zenodo packet contains three pairs of prepared-state
single-shot IQ arrays with approximately $2.94\times10^4$ records per class.
Projection onto the frozen class-centroid axis and classification at its
midpoint gives accuracies $0.9951$, $0.9956$ and $0.9980$, with pooled $d'$
values $7.32$, $6.34$ and $8.63$. These are reproducible standard readout
controls. They are integrated endpoint samples, not time-resolved trajectories,
and supply neither a process residual, holonomy nor body action. The public
Murch thermal-readout packet contains accumulated probabilities and histograms
rather than individual trajectories. Ficheux reports the strongest
architecture: time-resolved records, a frozen stochastic master equation and
independent terminal tomography on one carrier. Its raw records are available
only on request; public movies are not a reusable numerical archive. Under the
declared direct-public rule, zero located packets pass the six-part
process-residual gate. This closes the current public-data search at the
Tau-scoring level while retaining Wang as a pipeline control and Ficheux as the
leading acquisition target.

A separate public optical packet provides a useful same-experiment control
without repairing this identifiability failure. Wen et al. reconstructed
single-photon amplitudes over more than $1.4$ million paths and released the
figure tables under CC0 \citep{Wen2026PathData}. Across the 100 published
action bins, a weighted linear fit of path amplitude against classical action
gives slope $-0.0287\pm0.0417$, consistent with the equal-magnitude postulate.
The measured phase follows the reported action prediction with RMS residual
$0.0218$ in units of $\pi$. These are strong standard path-integral controls.
They are not a Tau score: path length and action are supplied as separate
aggregate tables, not as joint path-level $(M,S,A,\theta)$ records. Hence a
conditional novelty test such as $I(A;M\mid S)>0$ is not identifiable, and the
optical action cannot be renamed as $\mathcal A_X^{PC}$.

The raw camera matrices nevertheless allow one bounded follow-up diagnostic.
Using the published polarization differences, probability normalization and
the quadratic vertex of the independent Fig.~S1C sample, we reconstruct 17
radial step-kernel classes and enumerate the full $17^5=1{,}419{,}857$ path
family.  The phase reduction reproduces the Fig.~S1C sample to RMS
$0.0120\pi$ after one common phase offset.  Relative to scalar total standard
action, five frozen path descriptors add $\Delta R^2=0.1758$ to the phase
residual regression.  The generated paths, however, share the same 17 measured
kernel classes.  A permutation test at that kernel level gives $p=0.144$;
there is therefore no resolved morphology increment on the validated phase
leg.  The image-level amplitude reduction gives a larger apparent increment
but fails the independent Fig.~S1C amplitude check and is excluded.  This is a
standard factorized-kernel diagnostic, not evidence for an independent Tau
body action or a Tau-specific quantum signal.

The released Figs.~S2C--S2D permit a corrected amplitude-side reconstruction
without the failed image-level normalization. Their 17 normalized radial
means have only $0.863\%$ coefficient of variation, yet the corresponding path
products retain $\Delta R^2=0.4962$ beyond scalar total action. A permutation
of the 17 shared measured kernel classes gives $p=0.004$. Propagating the
histogram-derived standard errors through 9,999 draws under a constant-
amplitude null gives $p=0.0009$, and replacing any single radial class by its
action-trend prediction leaves $\Delta R^2\geq0.2936$. Thus the amplitude
structure is reproducible from the authors' published reduction and is not a
single-class artifact. It remains a standard optical-kernel morphology signal,
not a Tau endpoint: all paths reuse the same 17 measurements and the standard
optical action is not an independently reconstructed $\mathcal A_X^{PC}$.

To avoid interpreting this five-descriptor fit after the fact, we next freeze
one theory-motivated candidate before comparison with the controls. The
candidate is the Dirichlet energy of the path-increment field,
\begin{equation}
 M_{\rm curv}(\gamma)
 =\sum_k\left(\Delta x_{k+1}-\Delta x_k\right)^2.
 \label{eq:wen-curvature-candidate}
\end{equation}
This is only a conditional Tau descriptor: the body program selects a
primitive Dirichlet form, whereas applying it to the increment field requires
a separately typed, body-derived edge carrier that is not proved here. With
scalar total action retained as the baseline,
Eq.~\eqref{eq:wen-curvature-candidate} adds only
$\Delta R^2=0.000405$. A permutation of the 17 shared radial kernel classes
gives $p=0.543$, and the candidate ranks fourth among the five single-
descriptor additions. Ordinary path length alone gives
$\Delta R^2=0.4826$. Thus the robust published-amplitude result is dominated
by a standard accumulated-step descriptor; the frozen curvature candidate is
not detected. This negative result narrows the admissible Tau morphology
family and prevents the broad five-variable signal from being promoted to a
Tau-specific effect.

An independent aggregate check is supplied by the released Fig.~S2B table.
It contains 17 measured radial phase points with standard deviations and a
200-point theoretical curve. Linear interpolation of the theory onto the
measurement grid gives an RMS residual of $0.01233\pi$; all points lie within
one reported standard deviation, with maximum absolute standardized residual
$0.4873$. This strengthens the standard phase reconstruction but not the Tau
identifiability claim: Fig.~S2B contains no same-index complex amplitude,
recovery-process record or independently reconstructed body action. It
therefore remains ineligible for the primary Tau score.

This gives a positive but conditional reopening architecture. On one
programmable qubit--ancilla carrier, reconstruct the characteristic function
of work by Ramsey interferometry, independently tomograph withheld endpoints,
and close at least one Bargmann or Uhlmann loop. Freeze the action map and
$A_*$ before reading the endpoint and loop data. The resulting packet would
test both the finite magnitude identity and connector transport without
cross-platform stitching. The audit therefore replaces an open-ended source
search by one finite design target; it is a feasibility result, not evidence
that Nature realizes finite A8b.

The three legs are also mathematically compatible on one explicit carrier.
Consider the closed qubit family
\begin{equation}
 H(q)=\frac{\Delta_0(1+m\cos q)}{2}\,
 \mathbf n(q)\!\cdot\!\boldsymbol\sigma,
 \qquad
 \mathbf n(q)=(\sin\theta\cos q,\sin\theta\sin q,\cos\theta),
 \label{eq:joint-qubit-source}
\end{equation}
with Gibbs records
$\rho(q)=e^{-\beta H(q)}/\Tr e^{-\beta H(q)}$.  On this same path, a sudden
step has the standard work characteristic function
\begin{equation}
 G_{q\to q'}(u)=
 \Tr\!\left[e^{iuH(q')}e^{-iuH(q)}\rho(q)\right],
 \label{eq:joint-work-characteristic}
\end{equation}
tomography gives the Uhlmann chord
$x_Q=2(1-\sqrt{F(\rho(q),\rho(q'))})$, and discrete Uhlmann parallel transport
around $q\in[0,2\pi]$ gives a closed-loop phase.  For the frozen control
$(\beta,\Delta_0,m,\theta)=(1.3,1,0.35,1)$ on 65 points, the audit finds a
nonconstant $G(u)$, nonzero step chords and loop phase
$\Phi_U=0.13427$ radians; the final transport-unitarity defect is below
$8\times10^{-16}$.

\paragraph{Joint-carrier compatibility result.}
Work-characteristic interferometry, endpoint state geometry and Uhlmann loop
transport can therefore be defined consistently on one Hamiltonian family
without combining different carriers.  The construction does not prove
$x_A=x_Q$: Eq.~\eqref{eq:joint-work-characteristic} is a terminal
thermodynamic work statistic, and the independent map from this statistic to
$\mathcal A_X^{PC}/A_*$ remains absent.  Its role is to remove protocol
incompatibility as a blocker and isolate the remaining source-law question.

There is, however, a finite obstruction to using ordinary dissipated work as
that missing source law. For a sudden step between Gibbs states, direct
substitution of $\log\rho(q)=-\beta H(q)-\log Z(q)$ gives
\begin{equation}
 \beta\bigl(\langle W\rangle-\Delta F\bigr)
 =D\!\left(\rho(q)\Vert\rho(q')\right).
 \label{eq:dissipated-work-relative-entropy}
\end{equation}
Thus dissipated work induces relative-entropy/Bogoliubov--Kubo--Mori
geometry, whereas the finite A8b target uses the Uhlmann/Bures chord. These
geometries are not related by one universal constant. On the frozen loop
above, the stepwise ratio
$D(\rho(q)\Vert\rho(q'))/x_Q(q,q')$ ranges from $2.1171$ to $2.4882$.

\paragraph{Ordinary-work bridge no-go.}
Even on the joint carrier, no path-independent common gain converts ordinary
dissipated work into the finite Uhlmann chord. Work interferometry remains a
valuable independent control and may constrain a richer source map, but the
identification $\beta W_{\rm diss}=x_A=x_Q$ is rejected. A successful bridge
must instead descend the already defined once-counted body action, or derive
a nontrivial state-dependent transformation before endpoint data are read.

The inherited body construction suggests a different, noncircular work
observable. Let $z$ be a source-side occupied defect coordinate and let its
once-counted positive body action be
\begin{equation}
 \mathcal A_X^{PC}(z)
 =\frac12\langle z,K_A(z)z\rangle,
 \qquad K_A(z)>0.
 \label{eq:body-action-potential}
\end{equation}
Its conjugate body-work one-form is not the terminal qubit energy change but
\begin{equation}
 \alpha_B:=d\mathcal A_X^{PC},
 \qquad
 x_A(z)=\frac{1}{A_*}\int_{0}^{z}\alpha_B
       =\frac{\mathcal A_X^{PC}(z)}{A_*},
 \label{eq:body-work-transgression}
\end{equation}
where the zero-defect reference, physical path family and $A_*$ are frozen
before state tomography. Exactness makes the integral path independent on the
tested simply connected body chart; failure of closure is itself a rejection
of this minimal bridge and must not be repaired by endpoint fitting.

\paragraph{Conditional body-work bridge.}
An independent generalized-force reconstruction of $\alpha_B$, including its
closure and normalization ledger, supplies the missing action ordinate without
using $\rho$, fidelity or QGT. Equation~\eqref{eq:tau-discriminator} can then
be tested blindly against withheld endpoint tomography and loop holonomy. This
is an operational realization of the inherited once-only body-energy
transgression, not a proof that a presently published qubit experiment has
measured it.

\paragraph{Resonator-factorization no-go.}
A readout resonator is not automatically a second physical axis. Let $c$
denote frozen controls and suppose its reported response satisfies
\begin{equation}
 R_{\rm res}=F(\rho,c).
 \label{eq:resonator-factorization}
\end{equation}
Then $D R_{\rm res}=D_\rho F\,D\rho$, so
$\ker D\rho\subseteq\ker D(\rho,R_{\rm res})$ and stacking the resonator
record with state tomography cannot increase rank beyond the state map.
Equivalently, the response contains no conditional information about a hidden
action coordinate once $(\rho,c)$ is known. Thus dispersive cavity shifts used
to reconstruct the qubit state, and dynamical responses designed to reconstruct
the QGT, cannot also serve as an independent $x_A$ ordinate.

A finite seven-candidate audit separates two questions. Dispersive readout,
measurement-instrument backaction and QGT-response protocols do not supply an
independent action axis
\citep{Filipp2009DispersiveReadout,Hatridge2013Backaction,Rudinger2021QILGST,
Tan2019QGT,Yu2019QGT,Zhao2020GeometricGate}. Calorimetric work measurement is
the one audited architecture that can add an independent energy-exchange
record \citep{Pekola2013Calorimetry}, but the cited construction does not
provide a public same-carrier endpoint/holonomy packet or a frozen map from
ordinary energy exchange to $\mathcal A_X^{PC}/A_*$. Thus independent
measurement rank is necessary but not sufficient for an independent $x_A$.
The positive requirement is a separately calibrated source-port force,
impedance, admittance or calorimetric backaction whose response does not factor
through $\rho$ or the QGT, followed by a body-action map frozen before
tomography.

\paragraph{Generalized-force factorization no-go.}
The phrase ``generalized force'' does not by itself meet this requirement. For
a standard controlled Hamiltonian family $H(\lambda)$, the measured force is
\begin{equation}
 \mathcal M_\mu=-\operatorname{Tr}\!\left(\rho\,\partial_\mu H\right).
 \label{eq:standard-generalized-force}
\end{equation}
Once $\rho$, $H$ and the control calibration are frozen,
$\mathcal M_\mu=F_\mu(\rho,H,\partial H)$ and its Jacobian row lies in the span
of those inputs.  Slow-drive corrections can expose Berry curvature and other
quantum-geometric response coefficients; this makes generalized-force
experiments valuable independent \emph{instrumental} checks of quantum
geometry, but not an information-independent ordinate for $x_A$.  Using the
same response both to reconstruct the QGT and to test $x_A=x_Q$ would therefore
be circular \citep{Kolodrubetz2017Geometry,Schroer2014TopologicalTransition}.

The shortcut reopens only if a separately calibrated residual
\begin{equation}
 \xi_\mu:=\mathcal M_\mu+
 \operatorname{Tr}\!\left(\rho\,\partial_\mu H\right)
 \label{eq:source-port-force-residual}
\end{equation}
is reproducibly nonzero, survives conditioning on state, controls and
instrument response, and is mapped to the once-counted body action with its
unit frozen before endpoint tomography.  A nonzero $\xi_\mu$ would first
diagnose an incomplete standard source-port model; it would become a Tau test
only after that independent body-action bridge is established.  Thus the
standard generalized-force/QGT route is closed, while calorimetric or
mechanical source-port architectures remain conditionally open.

\paragraph{Bounded public source-port verdict.}
We finally audited the strongest public thermal and mechanical alternatives
rather than extending the candidate census indefinitely.  Gunyh\'o et al.
publish a large single-shot thermal-detector archive, but the bolometer absorbs
the state-dependent resonator readout power; it is a physically distinct
detector, not a conditionally independent source/body ordinate
\citep{Gunyho2024ThermalReadout}.  Pulsed optomechanical tomography reconstructs
mechanical phase space with calibrated optical pulses, while collective
backaction-evading measurements establish mechanical tomography and
force-sensing capability \citep{Vanner2013MechanicalTomography,
OckeloenKorppi2016BAE}.  Neither publication supplies a joint applied-force
record surviving conditioning on the reconstructed state, a finite
Uhlmann/QGT endpoint and a frozen map to $\mathcal A_X^{PC}/A_*$.  The finite
public-source-port audit therefore has zero eligible packets.  This closes the
present public-data-only acquisition class; Eq.~\eqref{eq:body-work-transgression}
remains a future protocol or an internal source-law target, not an observed Tau
signal.

\paragraph{Inherited CHDF body-work theorem.}
The internal source side is stronger than the public-data status. Let
$L:X_B\rightarrow E_S$ be the established visible body map and let $H_B>0$ be
the frozen CHDF body Hessian. Exact constrained elimination gives
\begin{equation}
 \mathcal A_Q(y)=\inf_{Lb=y}\frac12\langle b,H_Bb\rangle
 =\frac12\langle y,K_{\rm leaf}y\rangle,
 \qquad K_{\rm leaf}=(LH_B^{-1}L^*)^{-1}>0.
 \label{eq:chdf-quotient-action}
\end{equation}
Consequently the body-work form is not an additional ansatz:
\begin{equation}
 \alpha_B=d\mathcal A_Q
 =\langle K_{\rm leaf}y,dy\rangle,
 \qquad d\alpha_B=0,
 \qquad \int_{0}^{y}\alpha_B=\mathcal A_Q(y).
 \label{eq:chdf-exact-body-work}
\end{equation}
It inherits the same parent action unit and is path independent on every
simply connected quotient chart. The symbolic replay uses
$k_R=59/64-\sqrt3/12$ and $k_T=1+\sqrt3/24$ and verifies positivity, closure
and equality of straight and piecewise path integrals exactly.

This closes only the existence and normalization of the body-side one-form.
Under the adopted PRRC-P1/PJR-P1 constitutive continuation, the protected
body-relative role is carried into the post-body packet once, with no duplicate
action or free relative scale. That working-completion statement is not a
theorem that Nature selects PRRC-P1. More importantly, neither
Eq.~\eqref{eq:chdf-quotient-action} nor PRRC-P1 constructs the observer
amplitude connector required by finite A8b. The implication
$\alpha_B\Rightarrow x_A=x_Q$ therefore remains forbidden: the former is now
an inherited body theorem, while the latter is still the common-source
connector claim to be proved or tested.

\paragraph{Finite-range realization and non-selection theorem.}
The body action above is quadratic and therefore unbounded, whereas the
normalized Uhlmann chord
\begin{equation}
 x_Q=2\bigl(1-\sqrt{F(\rho_+,\rho_-)}\bigr)
 \label{eq:uhlmannchordbound}
\end{equation}
satisfies $0\leq x_Q\leq2$. Consequently a raw global identity $x_A=x_Q$
is impossible on the complete CHDF body. It can hold only on an occupied
bounded sector, or after a separately derived bounded reparametrization of
the action coordinate.

That bounded sector is nonempty and admits an exact construction. For
$0\leq x\leq2$, define
\begin{equation}
 \rho_\pm(x)=\frac12\left(I\pm r(x)\sigma_z\right),\qquad
 r(x)^2=x-\frac{x^2}{4}.
 \label{eq:a8bmixedrealization}
\end{equation}
The states are positive and normalized, their root fidelity is
$\sqrt{F}=\sqrt{1-r^2}=1-x/2$, and Eq.~\eqref{eq:uhlmannchordbound} gives
$x_Q=x$. Thus the CHDF action coordinate has an exact finite A8b realization
on $0\leq x_A\leq2$.

Existence is not physical selection. Composing the construction with a
depolarizing channel replaces $r(x)$ by $\gamma r(x)$ for any
$0\leq\gamma\leq1$. Every member is CPTP and preserves normalization, but
\begin{equation}
 x_Q^{(\gamma)}=2\left(1-\sqrt{1-\gamma^2r(x)^2}\right)
 \label{eq:a8bdepolarizingfamily}
\end{equation}
varies with $\gamma$. Positivity, normalization and CP descent therefore do
not force the A8b member $\gamma=1$. The complete-edge result of Paper III is
a sufficient internal realization theorem once the corresponding edge is
occupied; it does not prove that the physical base--seed parent selects that
edge. The remaining source statement is precisely the occupation law
$V_{\rm phys}=V_{r(x)}$, or an independently derived bounded replacement for
$x_A$. The range and counterfamily checks are reproduced by
\texttt{audit\_chdf\_uhlmann\_range\_and\_selection.py}.

\paragraph{Exact-recovery selector.}
The recovery branch removes the remaining freedom inside the depolarizing
counterfamily if it is imposed on the full occupied qubit operator system.
Let
\begin{equation}
 \mathcal C_\gamma(\rho)=\gamma\rho+(1-\gamma)I/2,
 \qquad 0\leq\gamma\leq1.
 \label{eq:depolarizingselector}
\end{equation}
For the two orthogonal eigenstates of $\sigma_z$, trace distance contracts
from one to $\gamma$. If a CPTP map $\mathcal R$ satisfied
$\mathcal R\mathcal C_\gamma=\mathrm{id}$ on the full qubit code, trace-distance
contractivity under $\mathcal R$ would imply
\begin{equation}
 1\leq D(\mathcal C_\gamma\rho_0,\mathcal C_\gamma\rho_1)=\gamma.
\end{equation}
Hence $\gamma=1$; conversely, $\gamma=1$ has identity recovery. Exact
full-code CPTP recoverability is therefore necessary and sufficient for the
A8b member within this family. This is a genuine selector theorem, but not yet
a Nature-selection theorem: the physical base--seed law has not been shown to
require exact recovery on the full occupied operator system. Restricted state
families and correctable subalgebras must not be silently promoted to that
premise. The finite check is reproduced by
\texttt{audit\_depolarizing\_recovery\_selector.py}.

\paragraph{Protected logical recovery versus global loss.}
The phrase ``full code'' in the selector theorem refers to the full occupied
logical ROOT operator system, not to every morphological-body realization
degree of freedom. This distinction is required by the Tau architecture,
whose observer descent is generically noninjective. Let
\begin{equation}
 \mathcal H_B=\mathcal H_L\otimes\mathcal H_G,\qquad
 \mathcal E_O=\operatorname{Tr}_G,qquad
 \iota(\rho_L)=\rho_L\otimes\tau_G,
 \label{eq:protectedfactor}
\end{equation}
where $L$ is the protected logical ROOT factor and $G$ contains sourced
realization multiplicity. Then
\begin{equation}
 \mathcal R(\sigma_L)=\sigma_L\otimes\tau_G,qquad
 \mathcal R\mathcal E_O\iota=\iota,
 \label{eq:protectedrecovery}
\end{equation}
although $\mathcal E_O$ is globally noninjective: states
$\rho_L\otimes g_0$ and $\rho_L\otimes g_1$ with $g_0\neq g_1$ have the same
observer output. Thus exact recovery of the logical algebra is compatible
with genuine observer-relative information loss in the body realization
factor.

More generally, for an isometric body--observer--complement realization,
exact recovery of the protected code is equivalent to the complementary
output carrying no information about the logical state. This is the precise
condition needed here: loss may retain realization or gauge information, but
must be blind to the ROOT coordinate used by A8b. The construction proves
compatibility and fixes the selector's scope. It does not prove that the
physical base--seed law places all ROOT information in the protected factor.
That source-factorization statement is now the remaining physical premise.
The finite logical/realization split is reproduced by
\texttt{audit\_protected\_logical\_recovery\_split.py}.

\paragraph{ROOT--commutant factorization.}
The logical/realization split need not be introduced as an additional tensor
postulate once one factor of the earlier Clifford packet has been physically
designated as the logical ROOT factor on the occupied $M_8(\mathbb C)$ block. Let
\begin{equation}
 \mathcal A_L=M_2(\mathbb C)\otimes I_4.
\end{equation}
Its commutant in $M_8(\mathbb C)$ is
\begin{equation}
 \mathcal A_G=\mathcal A_L'=I_2\otimes M_4(\mathbb C).
 \label{eq:rootcommutant}
\end{equation}
The algebras commute, have complex dimensions four and sixteen, and their
products span all $64$ complex dimensions of $M_8$. Hence a selected ROOT
$M_2$ factor canonically determines the realization factor and the
factorization $\mathcal H_B\simeq\mathcal H_L\otimes\mathcal H_G$ on this
block. Equation~\eqref{eq:protectedfactor} is therefore an internal
consequence of the selected tensor packet rather than a further arbitrary
enrichment.

This result does not remove the known tensor-selection boundary. Pointer
records alone admit entangling conjugacies, and the current narrow parent does
not select the required oriented Clifford flag. The remaining physical
premise is now sharply the selection and occupation of the ROOT factor/flag;
once supplied, its realization complement and protected-recovery split follow
automatically. The finite commutant and spanning checks are reproduced by
\texttt{audit\_root\_commutant\_}\allowbreak
\texttt{factorization.py}.

\paragraph{Factor-designation no-go.}
The rooted tensor packet does not itself perform that designation. Its root
stabilizer $S_3$ permutes the three opposite-edge pairs, and therefore the
three constructed $M_2$ factors, transitively. No individual factor projector
is fixed by this action; the invariant subspace of the three-dimensional
label representation is only the scalar line $(1,1,1)$. Thus OHC-P1 and the
rooted TC-CVP1 germ conditionally select an oriented but unordered factor
triple, not one distinguished logical ROOT factor.

Selecting one factor requires a typed $S_3$-breaking datum. The weakest
admissible form is a source-owned or irreducible observer--source access
functional whose three factor weights are not all equal and whose unique
extremum is frozen before terminal scoring. This datum may designate a
logical access context without changing the body or adding a fourth factor.
Its physical existence is not proved at this stage. Central chirality belongs
to a separate odd extension and is later shown to be invisible to the current
terminal. The permutation no-go is reproduced by
\texttt{audit\_root\_factor\_designation\_}\allowbreak
\texttt{no\_go.py}.

\paragraph{Extremal observer-access selector.}
The preceding no-go can be sharpened without adding a terminal fit. After the
three factor outputs have been source-typed into one common qubit codomain,
consider the convex access family
\begin{equation}
 \mathcal C_w=\sum_{k=1}^{3}w_k\operatorname{Tr}_{\widehat{k}},
 \qquad w_k\geq0,\qquad \sum_{k=1}^{3}w_k=1,
 \label{eq:weightedfactoraccess}
\end{equation}
where $\operatorname{Tr}_{\widehat{k}}$ retains factor $k$ and traces out the
other two. Encode a logical qubit in factor $j$ while placing the other
factors in the maximally mixed state,
\begin{equation}
 \iota_j(\rho)=\rho_j\otimes I/2\otimes I/2.
\end{equation}
Direct substitution gives
\begin{equation}
 \mathcal C_w\iota_j(\rho)
 =w_j\rho+(1-w_j)I/2.
 \label{eq:accessdepolarizing}
\end{equation}
Thus the effective depolarizing parameter on the $j$th logical factor is
$\gamma_j=w_j$. By the exact-recovery selector above, a CPTP left inverse on
the full logical qubit algebra exists if and only if $w_j=1$. Within
Eq.~\eqref{eq:weightedfactoraccess}, exact A8b recovery therefore selects an
extremal simplex vertex $w=e_j$, not merely a unique largest weight.

This result supplies two useful controls. Symmetric access
$w=(1/3,1/3,1/3)$ exactly designates no factor, while an interior vector with
a unique maximum still cannot support exact full-algebra recovery. Under an
$S_3$ relabeling, both $e_j$ and the retained factor transform together, so
different source-typed observer contexts may designate different factors
without contradiction or alteration of the common body. The result is a
conditional selector theorem, not by itself a derivation that Nature occupies
a vertex. Qubit-terminal minimality and regular-event incidence below derive
the corresponding on-code purity and contextual designation. The finite controls and
permutation-covariance check are reproduced by
\texttt{audit\_extremal\_observer\_access\_}\allowbreak
\texttt{selector.py}.

\paragraph{Qubit-terminal minimality theorem.}
The vertex condition is not peculiar to the weighted family once the terminal
type and recovery scope are fixed. Let the observer terminal be one full
$M_2(\mathbb C)$ algebra, and let the restriction of a CPTP access map
$\mathcal C$ to one occupied logical qubit code admit a CPTP left inverse
$\mathcal R$. Contractivity under both maps gives, for every pair of code
states,
\begin{equation}
 D(\rho,\sigma)
 =D(\mathcal R\mathcal C\rho,\mathcal R\mathcal C\sigma)
 \leq D(\mathcal C\rho,\mathcal C\sigma)
 \leq D(\rho,\sigma).
\end{equation}
Hence the restriction is a trace-distance isometry. A reversible CPTP map
between equal-dimensional full matrix algebras is unitary conjugation, so on
the occupied code
\begin{equation}
 \mathcal C(\rho)=U\rho U^\dagger.
 \label{eq:qubitterminalunitary}
\end{equation}
It follows that exact full-code recovery through a minimal qubit terminal
forces factor-pure access on that code; convex averaging or depolarization is
excluded. Moreover, preserving two complete qubit factors jointly would
require a four-dimensional code and cannot be reversibly represented in one
two-dimensional terminal.

This removes factor purity as an independent premise after exact recovery and
minimal terminal type have been adopted. It does not by itself select which
member of the unordered factor triple is occupied, prove that the physical
parent realizes exact recovery, or constrain off-code body information. Thus
the remaining continuous datum at this stage is factor \emph{designation},
not an additional gain or purity law. The dimension and scope audit is
reproduced by
\texttt{audit\_qubit\_terminal\_minimality\_}\allowbreak
\texttt{selector.py}.

\paragraph{Observer--source commutant designation.}
The remaining factor label can be supplied relationally rather than by a
preferred label in the common body. Let $\mathcal N_O$ be the algebra carried
by the complementary observer output, equivalently the realization/noise
algebra that is allowed to vary while the logical terminal is exactly
recoverable. If this algebra is maximal on the realization complement of
factor $j$, then
\begin{equation}
 \mathcal N_O=I_2\otimes M_4(\mathbb C),\qquad
 \mathcal A_{\mathrm{ROOT},O}=\mathcal N_O'=M_2(\mathbb C)\otimes I_4.
 \label{eq:observercommutantselector}
\end{equation}
The double-commutant relation returns $\mathcal N_O$ and makes the designation
unique for that typed context. Repeating the construction for each of the
three factors gives three contexts in one $S_3$ orbit. Thus observers may
designate different members of the same unordered body factorization while
their descriptions remain related by the body's permutation symmetry.

Equation~\eqref{eq:observercommutantselector} identifies what performs the
designation: the observer--source complementary algebra, not a new body
factor, fitted gain or globally preferred factor name. It is a conditional
source theorem. Exact recovery implies complementary blindness to logical
information, but this commutant theorem alone does not force the
\emph{maximal} factor-complement algebra assumed here. A smaller
noise algebra would have a larger commutant and would not uniquely select one
qubit factor. The finite commutant dimensions and the three-context orbit are
reproduced by
\texttt{audit\_observer\_noise\_commutant\_}\allowbreak
\texttt{selector.py}.

\paragraph{Canonical-expectation completion.}
Maximal complement occupation is not an additional free choice once a typed
factor embedding, complete body occupation and the normalized trace are
fixed. For
\begin{equation}
 M_8(\mathbb C)=M_2(\mathbb C)\otimes M_4(\mathbb C),
\end{equation}
the normalized-trace-preserving conditional expectation that is bimodular
over the designated logical algebra is uniquely
\begin{equation}
 \mathbb E_L=\operatorname{id}_2\otimes\tau_4,
 \label{eq:canonicalexpectation}
\end{equation}
where $\tau_4$ is the normalized trace on the realization factor. Equivalently,
the terminal channel is partial trace over that factor, up to the conventional
normalization and identification of the logical output. Its minimal Kraus
rank is four, so the complementary Hilbert space is four-dimensional and its
full observable algebra is $M_4(\mathbb C)$.

Thus the maximal complement in
Eq.~\eqref{eq:observercommutantselector} follows
from the already declared complete-body, trace-preserving and bimodular
descent conditions after the logical embedding has been designated. The
continuous source gap is reduced again: it is the typed observer--source
incidence selecting the $M_2$ embedding itself. Non-bimodular, nontracial or
incomplete-body maps lie outside this theorem and can have different
complements. The dimension ledger is reproduced by
\texttt{audit\_canonical\_expectation\_}\allowbreak
\texttt{complement.py}.

\paragraph{Regular-event incidence closes the factor label conditionally.}
The required typed incidence is not absent from the existing Tau completion.
On the adopted $E_1/\mathrm{EREM}$ observer branch, the earlier regular-event
theorem supplies a nonzero compact event cochain $j_{OS}^{\rm evt}$ for every
regular finite interaction. Resolve that cochain over the three factor
sectors of the occupied tensor packet and define
\begin{equation}
 w_k=
 \frac{\|P_kj_{OS}^{\rm evt}\|^2}
 {\sum_{\ell=1}^{3}\|P_\ell j_{OS}^{\rm evt}\|^2}.
 \label{eq:eventfactorweights}
\end{equation}
Regularity makes the denominator nonzero. The extremal-access and
qubit-minimality theorems then imply that an exact full-qubit A8b terminal is
possible only when $w=e_k$ for one unique $k$. That nonzero support identifies
the typed $M_2$ embedding, after which the canonical expectation fixes the
full complement.

No extra factor label or terminal fit has been introduced. A null event
selects no factor. A regular event with support on several factor sectors is
a valid observer interaction, but it cannot be promoted to an exact
full-qubit A8b terminal; it belongs to a noisy, approximate or richer-terminal
branch. Consequently the designation chain is closed conditionally on the
adopted $E_1/\mathrm{EREM}$ event completion and exact quantum recovery. The
minimal parent still does not force either physical occupation premise. The
odd chirality extension is outside this even-terminal selection. The null,
mixed and three axis controls are
reproduced by
\texttt{audit\_regular\_event\_factor\_}\allowbreak
\texttt{selection.py}.

\paragraph{Approximate-recovery stability.}
Exact recovery is the zero-error endpoint, not a requirement on every
physical interaction. Adopt the explicit convention
\begin{equation}
 \|\mathcal R\mathcal C-\operatorname{id}\|_\diamond\leq\eta.
 \label{eq:recoverydiamonderror}
\end{equation}
For two orthogonal logical states, each recovered state is then within trace
distance $\eta/2$ of its target. The triangle inequality gives recovered-pair
distance at least $1-\eta$, while contractivity under $\mathcal R$ bounds it
above by the pre-recovery distance. In the weighted access family that
distance is $w_j$, hence
\begin{equation}
 \boxed{
 w_j\geq1-\eta,
 \qquad
 \sum_{k\neq j}w_k\leq\eta.}
 \label{eq:approxfactorbound}
\end{equation}
Thus an experimentally bounded recovery error supplies a predeclared bound
on off-factor incidence. At five-percent diamond error, for example, the
selected factor carries at least $95\%$ of the normalized access weight.
Outside the weighted family the same argument bounds distinguishability
contraction but does not by itself define unique factor weights. The finite
error ledger is reproduced by
\texttt{audit\_approximate\_factor\_selection\_}\allowbreak
\texttt{bound.py}.

\paragraph{Central-chirality blindness of the present terminal.}
The central $\mathrm{Cl}_7$ sign must be separated from the quantum terminal
actually constructed here. Keep the six flag generators
$\gamma_1,\ldots,\gamma_6$ fixed and choose
$\gamma_7=\pm z_1z_2z_3$. The two extensions have opposite central Clifford
volume, but the first six generators already generate all of
$M_8(\mathbb C)$. Therefore the three factor algebras, their states and
effects, Born pairing, unitary transport, recovery maps and A8b geometry are
identical as objects of the common $M_8$ terminal algebra.

Consequently central chirality is not a blocker for the even quantum-readout
claim of Paper VI. It becomes physical only after a typed parity-odd,
orientation-sensitive or Weyl terminal couples to the odd extension. Such a
future terminal still requires a source-owned orientation lock, and the
present blindness result neither selects its sign nor proves chirality to be
gauge or absent. The two extensions and common rank-64 terminal algebra are
reproduced by
\texttt{audit\_central\_chirality\_terminal\_}\allowbreak
\texttt{blindness.py}.

\subsection{Predeclared readiness scoring}

The alternative-source census can be ranked for protocol readiness without
pretending to score a Tau signal. We froze a 13-point rubric before evaluating
the candidates: same carrier (2), direct work/action-like response (2),
endpoint geometry (2), loop holonomy (2), independent body-action bridge (3),
reusable public numerical data (1), and errors/recovery (1). Partial
capabilities receive one point only where the corresponding two-point item is
physically present but incomplete. A hard gate overrides the total: eligibility
for $\Delta_{AQ}$ scoring requires the same carrier, full endpoint and loop
legs, and the independent three-point body-action bridge.

\begin{table}[ht]
\centering
\small
\caption{Alternative-source readiness. Scores measure protocol completeness,
not evidence for Tau Core.}
\label{tab:alternative-readiness}
\begin{tabular}{@{}lccc@{}}
\toprule
Candidate & Score / 13 & Hard Tau gate & Principal missing leg\\
\midrule
Zhao Xmon geometric gate & 8 & fail & body action, reusable data\\
Cottet Maxwell demon & 7 & fail & holonomy, body action\\
Viyuela Uhlmann qubits & 7 & fail & direct work, body action\\
Batalh\~ao NMR work & 6 & fail & holonomy, body action\\
Ghosh NMR mixed phase & 6 & fail & direct work, body action\\
\bottomrule
\end{tabular}
\end{table}

The ranking selects the Xmon experiment as the closest integrated protocol
architecture, while the Cottet and Viyuela experiments provide the strongest
complementary work/tomography and holonomy realizations. The Tau-ready count
is nevertheless zero. The score may guide source acquisition; it may not be
converted into $\Delta_{AQ}$, a significance, or a detection claim.

